% ============================================================
% Section 3: Methodology
% LR (Nguyen Hieu An): 3.1 Dataset + 3.2 Experimental Setup
% MS (Vo Hieu Chuong):  3.3 Metrics + 3.4 Statistical Analysis Plan
% ============================================================

\section{Methodology}
\label{sec:method}

% ---------------------------------------------------------------------
% 3.1 Dataset — LR (Nguyen Hieu An)
% ---------------------------------------------------------------------
\subsection{Dataset}
\label{sec:dataset}

Our dataset contains 261 user stories in the Connextra template
(``As a [role], I want [feature], so that [benefit]''). We did not write
these stories ourselves. Instead, we reverse-engineered them from 261
expert-written \texttt{.feature} files taken from the production test
suites of three open-source projects: Sylius, Apache
Fineract, and Diaspora.\footnote{\url{https://github.com/Sylius/Sylius},
\url{https://github.com/apache/fineract},
\url{https://github.com/diaspora/diaspora}} We first crawled all Gherkin
feature files from the three repositories, discarded files that the
\texttt{behave} parser rejected, and selected 261 scenarios. We then
converted each scenario back into a Connextra user story with acceptance
criteria. Each user story is therefore paired with exactly one
expert-written Gherkin file, which serves as its ground truth. This
design avoids expert bias: the reference Gherkin was written by the
projects' own engineers for real production code, not by us for this
study.

We picked these three projects because they cover three unrelated
business domains: e-commerce, core banking, and social networking
(Table~\ref{tab:dataset}). A model that merely memorizes the
conventions of one domain cannot score well on the other two, which
strengthens the external validity of our results.

\begin{table}[t]
\caption{Dataset composition (N = 261 Connextra user stories).}
\label{tab:dataset}
\centering
\begin{tabular}{llr}
\toprule
\textbf{Project} & \textbf{Domain} & \textbf{User stories} \\
\midrule
Sylius           & E-commerce        & 100 \\
Apache Fineract  & Finance/Banking   & 100 \\
Diaspora         & Social Network    & 61  \\
\midrule
\textbf{Total}   &                   & \textbf{261} \\
\bottomrule
\end{tabular}
\end{table}

Listing~\ref{lst:example-us} shows one user story from the Sylius
subset, and Listing~\ref{lst:example-gt} shows the expert-written
Gherkin file it was derived from.

\begin{lstlisting}[float=ht, basicstyle=\ttfamily\scriptsize, frame=tb, captionpos=b, aboveskip=1em, belowskip=1em, breaklines=true, columns=fullflexible, caption={Example Connextra user story (id 1, Sylius, excerpt).},label={lst:example-us}]
As a store owner,
I want the checkout process to prevent
customers from completing an order with a
shipping method that has become unavailable,
So that customers cannot place orders using
shipping methods that are no longer valid or
offered.

Acceptance Criteria:
- When a customer has selected a shipping
  method and payment method during checkout,
  but that shipping method is subsequently
  disabled entirely, the system should block
  order confirmation [...]
\end{lstlisting}

\begin{lstlisting}[float=ht, basicstyle=\ttfamily\scriptsize, frame=tb, captionpos=b, aboveskip=1em, belowskip=1em, breaklines=true, columns=fullflexible, caption={Corresponding expert-written Gherkin ground truth (excerpt).},label={lst:example-gt}]
@checkout
Feature: Preventing placing an order with a
    disabled shipping method

  Background:
    Given the store operates on a single
      channel in the "United States" named
      "US Web Store"
    And the store has a product "Ubi T-Shirt"
      priced at "$19.99"
    And the store has "Raven Post" shipping
      method with "$4.00" fee
    And the store allows paying "Offline"

  @api @ui
  Scenario: Being prevented from placing an
      order with a shipping method that's
      disabled after completing the shipping
      method choice step
    Given I added product "Ubi T-Shirt" to
      the cart
    And I addressed the cart
    And I chose "Raven Post" shipping method
      and "Offline" payment method
    But this shipping method has been disabled
    When I try to confirm my order
    Then I should not be able to confirm order
      because the "Raven Post" shipping
      method is not available
  [...]
\end{lstlisting}

% ---------------------------------------------------------------------
% 3.2 Experimental Setup — LR (Nguyen Hieu An)
% ---------------------------------------------------------------------
\subsection{Experimental Setup and Pipeline}
\label{sec:pipeline}

All generations use \texttt{gpt-5.4-mini}, pinned to the snapshot
\texttt{2026-03-17}, through the official OpenAI Chat Completions API.
We set \texttt{temperature} to 0 so that decoding is deterministic: the
same input always produces the same output, which makes every run of
our pipeline reproducible. We impose no explicit max-token cap, so the
model is never forced to truncate a scenario mid-file. Pinning the
snapshot protects the study against silent model updates during the
experiment window.

We prompt the model with domain-specific few-shot prompting. Each of
the three domains gets (a) its own system message, which places the
model in the role of a QA engineer on that specific project, and (b)
exactly one example pair---a user story and its expert-written
Gherkin---taken from that same project. We chose this strategy because
the three projects write Gherkin in very different styles: Fineract
relies on dense data tables, Diaspora writes concrete UI steps with
CSS selectors in the Capybara style, and Sylius uses the
business-level language of Behat. A single generic example cannot
teach the model these project conventions, and prior work reports that
few-shot examples sharply reduce syntax errors in generated BDD
scenarios~\cite{karpurapu2024bddautomation}. Listing~\ref{lst:prompt} shows the
prompt template verbatim; the placeholders in braces are filled per
domain at run time.

\begin{lstlisting}[float=ht, basicstyle=\ttfamily\scriptsize, frame=tb, captionpos=b, aboveskip=1em, belowskip=1em, breaklines=true, columns=fullflexible, caption={Prompt template (verbatim). Placeholders in braces are instantiated per domain.},label={lst:prompt}]
System: {domain_specific_system_message}

User:
Generate a BDD Gherkin file for the following
user story.
Follow the exact style, wording, and
conventions shown in the Example.

EXAMPLE USER STORY:
{example_user_story_from_same_domain}

EXAMPLE GHERKIN:
{example_gherkin_from_same_domain}

---
NOW GENERATE FOR THIS USER STORY:
{user_story_text}
\end{lstlisting}

The pipeline runs in five steps. First, we load the 261 user stories
from \texttt{data/full\_user\_stories.csv}. Second, for each story we
build the domain-specific prompt and call the API. Third, we strip the
Markdown code fence (\texttt{```gherkin}) when the model wraps its
output in one; we make no other change to the generated text. Fourth,
we store every generated file in
\texttt{results/full\_gpt54mini\_output.csv} and append one log record
per API call---timestamp, model version, token counts, and cost---to
\texttt{results/full\_api\_log.jsonl}. Fifth, we score each generated
file on the two metrics defined in Section~\ref{sec:metrics}:
executability via \texttt{behave --dry-run} and semantic similarity
via \texttt{all-MiniLM-L6-v2} embeddings.

The pipeline also handles transient API failures. When a call fails or
times out, the script waits 60 seconds and retries the same story
until it succeeds, so no sample is dropped and no result is ever
filled in by hand. A one-second pause between calls keeps us under the
API rate limit, and the output file is saved after every sample, so a
crash loses no data. In the full run, all 261 calls returned a valid
response (N = 261; nothing was excluded).

Our replication package, including the dataset, prompts, scripts, and
raw results, is publicly available at:
\url{https://github.com/simonhoang611/RT-SWT-004-Group4}.

% ============================================================
% 3.3 Metrics — MS (Vo Hieu Chuong)
% ============================================================
\subsection{Metrics}
\label{sec:metrics}

We evaluate the quality of LLM-generated Gherkin scenarios using two complementary metrics,
each targeting a distinct research question.

% ----------------------------------------------------------
\subsubsection{Metric 1: Skeleton Cosine Similarity (RQ1 — Semantic Quality)}
\label{sec:metric-skeleton}

\textbf{What it is.}
Skeleton Cosine Similarity measures the semantic similarity between an LLM-generated
Gherkin scenario and its reference scenario by comparing their \emph{dense vector
representations} (embeddings).
We use the \texttt{all-MiniLM-L6-v2} sentence-transformer model~\cite{reimers2019sentencebert}
to encode each scenario into a 384-dimensional vector, then compute the cosine similarity
between the two vectors.
A score of 1.0 indicates identical meaning; 0.0 indicates orthogonal semantics.
We adopt a threshold of $\tau = 0.85$, consistent with prior work on semantic adequacy
in test-case generation~\cite{fraser2011evosuite}.

\textbf{Why the Skeleton variant.}
Raw Gherkin text contains structural keywords (\texttt{Feature:}, \texttt{Scenario:},
\texttt{Given}, \texttt{When}, \texttt{Then}, \texttt{And}) that are syntactically
mandatory but carry no domain-specific meaning.
Including these keywords biases the embedding toward formatting rather than business logic.
To isolate the \emph{core business logic}, we strip all Gherkin keywords before encoding,
producing a \emph{skeleton} string that contains only the natural-language step descriptions.

Table~\ref{tab:skeleton-example} illustrates this process with a concrete example.

\begin{table}[h]
\centering
\caption{Skeleton extraction example. Gherkin keywords (bold) are removed;
only the natural-language content is retained for embedding.}
\label{tab:skeleton-example}
\renewcommand{\arraystretch}{1.3}
\begin{tabular}{p{0.28\linewidth} p{0.28\linewidth} p{0.28\linewidth}}
\toprule
\textbf{Original Gherkin Step} & \textbf{Skeleton (after stripping)} & \textbf{Embedding (illustration)} \\
\midrule
\texttt{\textbf{Given} the user is on the login page}
  & \texttt{the user is on the login page}
  & $[0.23,\,-0.11,\,\ldots,\,0.87]$ \\
\addlinespace
\texttt{\textbf{When} the user enters invalid credentials}
  & \texttt{the user enters invalid credentials}
  & $[0.31,\,0.05,\,\ldots,\,-0.42]$ \\
\addlinespace
\texttt{\textbf{Then} an error message is displayed}
  & \texttt{an error message is displayed}
  & $[0.19,\,-0.27,\,\ldots,\,0.61]$ \\
\bottomrule
\end{tabular}
\end{table}

The final Skeleton Cosine Similarity score for a scenario is the cosine similarity between
the mean-pooled embedding of all skeleton steps of the generated scenario and that of the
reference scenario.

\textbf{Formal definition.}
Let $\mathbf{g}$ and $\mathbf{r}$ be the mean-pooled embeddings of the generated and
reference skeletons, respectively.
The Skeleton Cosine Similarity is defined as:
\begin{equation}
  \text{SkelCos}(\mathbf{g}, \mathbf{r})
    = \frac{\mathbf{g} \cdot \mathbf{r}}{\|\mathbf{g}\| \cdot \|\mathbf{r}\|}
\label{eq:skelcos}
\end{equation}
We evaluate $H_0^{\text{RQ1}}: \tilde{\mu} \leq 0.85$ (the population median does not
exceed the threshold) against $H_1^{\text{RQ1}}: \tilde{\mu} > 0.85$.

% ----------------------------------------------------------
\subsubsection{Metric 2: Executable Syntax Rate (RQ2 — Syntactic Correctness)}
\label{sec:metric-esr}

\textbf{What it is.}
The Executable Syntax Rate (ESR) is the proportion of generated Gherkin scenarios that
parse and load without errors in a real test runner.
Formally:
\begin{equation}
  \text{ESR} = \frac{|\{s \in S \mid \texttt{behave --dry-run}(s) = \textsc{pass}\}|}{|S|}
\label{eq:esr}
\end{equation}
where $S$ is the set of all generated scenarios.

\textbf{Why \texttt{behave --dry-run} instead of a static linter.}
Static linters such as \texttt{gherkin-lint} validate only the superficial token structure
of a \texttt{.feature} file.
They may pass files that contain step-definition references that are undefined, ambiguous,
or semantically malformed, which would cause actual test execution to fail.
By contrast, \texttt{behave --dry-run} performs a full parse-and-load cycle: it checks
Gherkin syntax, resolves step definitions against the step-definition registry, and reports
any \texttt{MissingStepImplementationError} or \texttt{AmbiguousStep} at dry-run time.
This gives a \emph{runtime-faithful} measure of syntactic correctness.

We set the acceptability threshold at $\tau_{\text{ESR}} = 0.80$ (i.e., $\geq$80\% of
scenarios must be executable), following industry guidelines for minimum viable test
coverage~\cite{wohlin2012experimentation}.
We evaluate $H_0^{\text{RQ2}}: p_{\text{ESR}} \leq 0.80$ against
$H_1^{\text{RQ2}}: p_{\text{ESR}} > 0.80$.

% ============================================================
% 3.4 Statistical Analysis Plan — MS (Vo Hieu Chuong)
% ============================================================
\subsection{Statistical Analysis Plan}
\label{sec:stats}

We use a significance level of $\alpha = 0.05$ for all hypothesis tests.
Two separate tests are applied, one per research question, selected based on the
distributional characteristics of each outcome variable.

% ----------------------------------------------------------
\subsubsection{RQ1 — Wilcoxon Signed-Rank Test (one-sample)}

We test whether the \emph{population median} of Skeleton Cosine Similarity exceeds
the threshold $\mu_0 = 0.85$.

\textbf{Why Wilcoxon (not a $t$-test).}
Cosine similarity scores are bounded in $[-1, 1]$ and, in practice, tend to cluster
near the upper end of the range.
Preliminary inspection of our data (see Figure~\ref{fig:rq1-dist}) reveals a
left-skewed distribution, violating the normality assumption of the one-sample $t$-test.
The \emph{Wilcoxon signed-rank test} is a non-parametric alternative that compares
medians without requiring a normal distribution~\cite{arcuri2011practical}.
It ranks the absolute deviations from $\mu_0$ and tests whether the sum of positive
ranks significantly exceeds that of negative ranks.

\textbf{Hypotheses (RQ1):}
\begin{align}
  H_0^{\text{RQ1}} &: \tilde{\mu}_{\text{SkelCos}} \leq 0.85 \notag \\
  H_1^{\text{RQ1}} &: \tilde{\mu}_{\text{SkelCos}} > 0.85 \notag
\end{align}
We report the one-sided $p$-value and Cliff's $\delta$ as the effect size
($|\delta| < 0.147$: negligible; $0.147$--$0.33$: small; $0.33$--$0.474$: medium;
$\geq 0.474$: large~\cite{romano2006exploring}).

% ----------------------------------------------------------
\subsubsection{RQ2 — Exact Binomial Test (one-sample)}

We test whether the true proportion of executable scenarios exceeds the threshold
$p_0 = 0.80$.

\textbf{Why Exact Binomial.}
Each scenario is independently assessed by \texttt{behave --dry-run} and yields a
binary outcome: \textsc{pass} or \textsc{fail}.
This is a Bernoulli trial with unknown success probability $p$.
The Exact Binomial test is the standard choice for this data type: it computes the
exact probability of observing at least $k$ successes out of $N$ trials under
$H_0: p = p_0$, without relying on asymptotic (normal) approximations that may be
unreliable when $p$ is close to 0 or 1~\cite{agresti2002categorical}.

\textbf{Hypotheses (RQ2):}
\begin{align}
  H_0^{\text{RQ2}} &: p_{\text{ESR}} \leq 0.80 \notag \\
  H_1^{\text{RQ2}} &: p_{\text{ESR}} > 0.80 \notag
\end{align}
We report the one-sided $p$-value and the 95\% Clopper-Pearson confidence
interval~\cite{clopper1934use} for the true proportion.
